\section{Empirical result}\label{Empirical result}

Due to the one-dimensional structure of MPS, Born machines are particularly well-suited for sequence modeling. Among various types of sequential data, the biological gene sequence is one of the most vital and holds immense value for humans. However, it is challenging to obtain experimentally and often incurs significant costs to extract limited data. Therefore, using generative models to discern the underlying probability distribution in gene sequences and generate new samples can enhance the efficiency of exploring new gene sequences and reduce costs. In this study, we will evaluate the performance of Born machines with trainable POVMs for the unsupervised learning of RNA sequences. We uses the bactaria's 5S rRNA with average length 120 and total vocabulary size of four(\textit{A}: Adenine, \textit{T}: Thymine, \textit{C}: Cytosine, \textit{G}: Guanine) from open source database \textit{5SrRNAdb}\cite{szymanski20165srnadb}. The training set and test set are unbiasedly divided from the dataset.

\begin{figure}[htbp]
\begin{center}
\includegraphics[width=250pt]{samples.pdf}
\caption{In this figure, each horizontal line represents an RNA sequence, displayed vertically. The four colors represent four types of nucleotides. (a) Generated samples by Born machine with POVM embedding. (b) Generated samples by GPT-2. (c) Real data from the test set.}
\label{fig: generation}
\end{center}
\end{figure}

\figref{fig: generation} generated by the Born machine with POVMs and GPT-2 appear qualitatively similar to the real test set, indicating both models capture visible patterns in the sequence data. Given this similarity, we proceed to evaluate each model’s performance using quantitative metrics such as NLL, single-site probabilities, and Pearson correlations. \figref{fig: loss} demonstrates that increasing the physical dimension in POVMs reduces the converged NLL values, even with a fixed boundary dimension. Furthermore, it reveals that higher physical dimensions extend the upper limit of bond dimension, enabling further NLL minimization.

\begin{figure}[htbp]
\begin{center}
\includegraphics[width=200pt]{loss.pdf}
\caption{The NLL objective function is minimized using \textit{Adam} optimizer with adjustable learning rate for bond/physical dimension separately. Result shows that the converged NLL is decreased by jointly increasing the bond dimension and the physical dimension. The green dashed line represents the converged NLL of GPT-2 model with similar parameter size.}
\label{fig: loss}
\end{center}
\end{figure}

One of the features of RNA data distribution is that, the marginalized joint probability distribution among various sites reveals the underlying contact prediction information\cite{trinquier2021efficient}. Contact prediction in RNA refers to the process of determining which residues or nucleotides in an RNA molecule are in close proximity to each other in the molecule's three-dimensional structure. To delve deeper into evaluating the performance of the Born machine with trainable POVMs and discovering contact prediction, one can leverage its efficient marginalization ability \eqnref{eq: single site} and then compare it to statistical frequencies extracted from the dataset. In particular, we will study two \textit{statistical features}:
\begin{itemize}
\item the single-site probability distribution $p(x_i)$,
\item and the two-site Pearson correlation
\begin{equation}\label{eq: corre}
    c(x_i,x_j)=\log\frac{p(x_i,x_j)}{p(x_i)p(x_j)}.
\end{equation}
\end{itemize}
Each sub-figure in \figref{fig: corre} plots the model-predicted statistical feature (vertical axis) against the data-indicated statistical feature (horizontal axis). \figref{fig: corre}(a-e) is for $p(x_i)$ where each point in the plot corresponds to a different choice of site $i$ and a token $x_i$ on that site. \figref{fig: corre}(f-j) is for $c(x_i,x_j)$ where each point in the plot corresponds to a choice of pair $(i,j)$ and tokens $(x_i,x_j)$ on that pair. If the model learns all statistical features from data perfectly, we should expect all points to fall along the diagonal dashed line (on which the model prediction coincides with the data statistics).  The result in \figref{fig: corre} indicates that as the physical dimension increases, both the single-site probability and correlations align more closely with the test set. This indicates that the Born machine equipped with trainable POVM embedding will be more powerful in modeling token correlations across different sites in the training data, compared to the fixed one-hot embedding approach.


We also benchmarked our quantum model against the classical sequential generative model, the Generative Pretrained Transformer-2 (GPT-2), with a comparable parameter size. The GPT-2 configuration is initialized from the \textit{Hugging Face transformer} library\cite{gpt2lmhead} with $d_{\text{embedding}} = 72$, $n_{\text{layer}} = 12$, $n_{\text{head}} = 4$, which contains approximately 0.8 million parameters in total, and is integrated with the GPT-2LMHead model. For the Born machine, the total number of parameters can be estimated as the product of the bond dimension, physical dimension, and sequence length, i.e., $d_{\text{bond}}^2 \times d_{\text{phys}} \times L$. For example, our largest configuration uses $d_{\text{bond}} = 20$, $d_{\text{phys}} = 16$, and $L = 120$, leading to approximately 0.77 million parameters—comparable to GPT-2. In \figref{fig: loss} and \figref{fig: corre}, the Born machine achieves lower NLL and comparable performance in approximating single-site probabilities relative to GPT-2 under our current settings. However, it shows weaker proficiency in correlation approximation. Attention-based networks like GPT-2 are designed with an all-to-all connectivity structure without inherent dimensionality constraints, making them particularly effective at capturing long-range correlations.

\begin{figure}[htbp]
\begin{center}
\includegraphics[width=240pt]{corre.pdf}
\caption{Compare (a-e) the single-site probability $p(x_i)$ (upper row) and (f-j) the two-site correlation $c(x_i,y_j)$ (lower row) of the model with those from the dataset. GPT-2: (a)\&(f). One-hot embedding (baseline): (b)\&(g). Trainable POVM embedding (ours) with the physical dimension: (c)\&(h) $p=4$, (d)\&(i) $p=8$, (e)\&(j) $p=16$. The dashed line represents the $y=x$ reference line.}
\label{fig: corre}
\end{center}
\end{figure}


\section{Summary}

Generative models are a vital branch of machine learning that strive to understand the inherent probability distribution of data. Unlike discriminative models, they focus on capturing the data's intrinsic structure. This work continues on studying quantum-inspired model called the Born machines for generative modeling. These machines use the quantum wave function to model underlying data distribution over the matrix product state (MPS) framework. The main innovation in this study is the extension of the token embedding method. Instead of embedding each token directly by a corresponding tensor index, this work introduces the concept of trainable positive operator-valued measurements (POVMs). By combining MPS with trainable embeddings, Born machines can achieve better performance and extract deeper correlations from datasets. The source code and raw data are available in the GitHub repository\cite{hou2023github}.

Our new Born machine architecture leverages quantum wave functions as its backbone, with tokens encoded as quantum measurement matrices. The joint probability of each data point is expressed as the expectation value of a quantum measurement. A key feature of the Born machine is its capability for autoregressive generation, enabling sequence sampling in arbitrary order—making it versatile for tasks such as masked token sampling and anomaly detection. While our current implementation is classically simulable via matrix product states (MPS), the true advantage may arise when sampling is carried out on a real quantum device. A promising strategy is to train the model efficiently on classical hardware using MPS, and then transfer the trained state to a quantum computer via an equivalent quantum circuit. As demonstrated in \cite{2022PhRvX..12a1047H, 2024QS&T....9a5012R}, each MPS tensor can be mapped to local quantum gates with ancilla insertions and measurements in a quantum circuit, representing the MPS as a real quantum state on the quantum computer. Furthermore, the adaptive POVM method introduced in this work can be implemented via generalized measurements on an extended Hilbert space \cite{preskill1998lecture}, enabling efficient sampling and potentially leveraging quantum advantage. Since POVMs can be directly realized in practice, this approach can be extended to the Quantum Circuit Born Machine (QCBM) architecture for implementation on quantum hardware \cite{PhysRevA.98.062324, 2019npjQI...5...45B}.

In terms of application, our study focuses on the unsupervised learning of RNA sequences. Results indicate that by increasing the physical dimension in the POVMs, the Born machine can achieve better performance, capturing deeper correlations in the RNA data. However, there are still unresolved issues within the current framework of the Born machine. For instance, its present structure cannot directly handle data with significant length variances. Moreover, because the current Born machine is constructed using a one-dimensional MPS with a discretized embedding method, it's awkward to directly extend it to handle higher-dimensional and continuous data types. Additionally, there's the challenge of transferring the well-trained MPS and corresponding POVMs to a quantum device for efficient sampling.

\begin{acknowledgments}
We acknowledge the helpful discussions with Rose Yu, Tian-Xiao Hu and Zhao-Yi Zeng. 
We acknowledge the OpenAI GPT4 model for providing editing suggestions throughout the process of writing this paper. The research is supported by the NSF Grant No.~DMR-2238360.
\end{acknowledgments}


